\scanpage{23}\mspage{22}
Il en résulte que $\overline{\{\deleted{\unsure{b}}a\}}=\overline{\{b\}}\Longleftrightarrow a=b$ (donc l'espace top $\mathcal L$ satisfait l'axiome $T_0$, il est ``présobre'' i.e. s'injecte dans le présobrifié) -- je ne sais si c'est
le cas pour tout tronçon, ni si un tronçon, avec sa top. canonique, est sobre (donc, de plus, tte partie fermée \uline{irréductible} aurait un pt générique \dots)
\figure{\begin{tikzpicture}[x=1cm,y=1cm]\draw (0,0)--(1.25,-.03)--(1.56,1)--(3.4,1);\draw (1.18,-.01)--(1.49,1);\end{tikzpicture}}
\margin{\unsure{A|À} voir,}
les ``segment'' $[a,b]$ ne sont pas néc. fermés pour la top. canonique$^{\dagger}$\unsure{,|.}
Ce n'est que lorsque $\mathcal L$ est \uline{tot} ordonné qu'on peut conclure que les pts sont fermés (réciproque ?) et en fait, $\mathcal L$ est $\unsure{=|\simeq}$ \uline{séparé}, et de plus les $[a,b]$, $\mathcal L_{\geq a}$, $\mathcal L_{\leq\deleted{\unsure{b}}a}$ sont fermés.

On peut définir une autre topologie, ayant comme base fondamentale de fermés les $[a,b]$, $\mathcal L_{\geq a}$, $\mathcal L_{\leq a}$ (dont on prend les intersections quelc. de réunions finies, pour avoir tous les fermés)

% pale vertical watermark
\rotatebox{90}{\textsc{\unsure{ng Extra Stro}}}
